\section{Overview}\label{sec:overview}
% NEW SECTION 2 (restructure task T3, 2026-07-24) per the PI's restructuring email
% + the 2026-07-24 restructure meeting [31:30]: definitions -> problem definition -> challenges.
% The definitions and the problem paragraph MOVED here verbatim-in-substance from
% 03_methodology.tex's retired \subsection{Problem Formulation}; every formula,
% number, and % src: comment travelled with them. Two content-level changes only:
% (1) the beats are re-ordered into "data and representation" then "the service
% allocation the data induces" (the email's shape), which splits the original's
% single-sentence pairing of the grid + active-unit definitions and moves the
% "D and S are deterministic aggregations" clause next to where D and S are
% defined (its duplicate gloss "(demand counts a unit's recorded pickups, ...)"
% dropped as an immediate repeat); (2) \emph{} on defined terms upgraded to
% \textbf{} at first definition (cGAIL style), plus one added sentence defining
% \textbf{collective} fairness, the endorsed term the challenges below lean on
% (ALGORITHM_FACTS: "fairness is collective (corpus-level)"). A commented-out
% earlier variant of the corpus paragraph was not carried over (it is in git
% history at d88f1a9:paper/sections/03_methodology.tex, lines 38-50).
% The label sec:problem MOVED here with the content so the appendix's three
% "\S\ref{sec:problem}" pointers keep resolving; they now read as this section.
\label{sec:problem}

\textbf{Data and representation.} FATE edits human-generated spatial-temporal
data: here, a \textbf{corpus} $\mathcal{T} = \{\tau\}$ of real
% (HSTD acronym is defined in the abstract and §1; spelled out here, not re-defined)
per-driver taxi trajectories, recorded on a city discretized into grid cells
$g$ and hourly time blocks $t$ (see Appendix~\ref{app:editor} for cell
resolution and per-city grid dimensions). A \textbf{trajectory}
$\tau = (s_1, \dots, s_T)$ is a time-ordered sequence of states: each
\textbf{state} $s_j = (g_j, t_j)$ records the cell and time block its driver
occupied while seeking a passenger, and the final state $s_T$ is the
\textbf{pickup}.
% src: PAPER/argument/02_datasets.md (0.01° grid, hourly blocks)
The imitation setting is supervised (\S\ref{sec:downstream}): there is no
reward function, and the recorded state transitions are the only training
signal. A policy $\pi$ is trained by behavior cloning to
predict a driver's next state from their recent states, so the policy can only
be as fair as the corpus it imitates. The trajectories are the editable
objects.

\textbf{The service allocation the data induces.} A pair $i=(g,t)$ with
recorded activity is an \textbf{active unit}, and $N$ denotes their count
($N \approx 34{,}500$ on Shenzhen, the primary city). Each active unit carries a
\textbf{demand} $D_i$ and a \textbf{supply} $S_i$: demand is its recorded
pickups, and supply is active-taxi presence aggregated over the
$5{\times}5$ cell neighborhood, matching how a cruising taxi is visible to
nearby demand. The \textbf{service ratio}
$Y_i = S_i / \max(D_i, d_0)$ is service received relative to demand, where
the demand floor $d_0 = 0.5$ prevents near-empty cells from producing
unbounded ratios. Both $D$ and $S$
are deterministic aggregations of $\mathcal{T}$, so moving where a driver
seeks and picks up moves demand and supply with it.
% src: PAPER/supply-lift/LIFT_ALGORITHM_REFERENCE.md §3 (floors, 5×5 presence); §4.1 (N = 34,524 active units, Shenzhen)

Each cell additionally carries demographic covariates $x_i$ describing its
neighborhood: housing price, per-capita compensation, and migrant population
share. These covariates resolve at district granularity, which bounds the
resolution of any demographic analysis built on them (\S\ref{sec:objective}).
Fairness throughout this paper is \textbf{collective}: a property of the
service allocation the whole corpus induces over units and demographic
groups, never a property of one trajectory.
% src: PAPER/argument/02_datasets.md

\textbf{Problem definition.} Given a corpus $\mathcal{T}$ whose service
allocation is demographically inequitable, produce an edited corpus
$\mathcal{T}'$ that (i) is measurably fairer; (ii) preserves fidelity at the
level of individual trajectories: an edited trajectory must still be
recognized as its original driver by a frozen driver-identity
discriminator; and (iii) transfers its fairness into a policy trained on it by
imitation. FATE is a \textbf{data-augmentation} method: it edits a small,
attribution-targeted slice of the real trajectories and generates nothing. At
most $k$ trajectories are edited ($k = 10{,}000$ on the primary city), each
edit bounded to an $L_\infty$ ball of radius $\varepsilon = 2$ grid cells.
% 2026-07-24 (T9, consistency micro-sweep): "each perturbed within" -> "each edit
% bounded to". FATE's own edits are never called perturbations or attacks (that
% vocabulary is reserved for the adversarial lineage and the baseline arms); the
% replacement matches C2's "every edit bounded to a small neighborhood" and Fig.~2's
% "every edit stays within eps cells".
% src: LIFT_ALGORITHM_REFERENCE.md §3 (EPSILON_BALL, budget)
The edited slice is then upweighted during downstream training so that its
fairness is not averaged away (\S\ref{sec:downstream}).

\textbf{Challenges.} Editing such a corpus into a fairer one is hard for five
concrete reasons; each dictates one component of the method
(\S\ref{sec:method}), and we refer back to them by label throughout.
% Stacked in running text, deliberately NOT an itemize environment
% (the 2026-07-24 restructure meeting [32:12]: "we don't need to kind of put them in this itemized
% LaTeX environment. We can just stack them all together"). Five challenges,
% 1:1 with the methodology blocks, per the adjudicated mapping in
% restructure/MEETING_DIGEST.md ADJ-3. Numbering keeps the old intro list's
% C2-C5 identities and rewrites C1 from "the data is scarce" to the budget
% challenge (the meeting's "unique challenge", [06:02]).
% 2026-07-24 (T10, page-budget lever 3): each challenge compressed to two tight
% sentences. Every load-bearing clause survives: C1's budget $k \ll |T|$ and the
% whole-corpus difficulty; C2's read-as-its-driver requirement, the bound, and
% the discriminator scoring INSIDE the objective (never a gate); C3's
% demand-unexplained variation; C4's empirical "a demand-only editor does exactly
% that on our data" plus the adds-real-presence remedy; C5's averaged-away
% mechanism plus tracing the recovery to the edits rather than the reweighting.
% Dropped as recoverable elsewhere: C1's "discarding the biased slice or
% regenerating the corpus discards the asset" (the data-augmentation framing of
% the problem definition above), C2's "no longer resembles real driving" gloss on
% the sentence that follows it, C4's "redistributing service among better-served
% areas" (stated in full in \S\ref{sec:leveling}), C5's "never reaches the
% policy" (the same fact as being averaged away). Stacked format and the lead-in
% sentence kept.
\textbf{C1 (budget).} Real demonstrations are expensive and irreplaceable, so
only a small share of the corpus may change. The difficulty is where to spend
a budget of $k \ll |\mathcal{T}|$ edits so that a few trajectories move the
fairness of the whole corpus (\S\ref{sec:editor}).
\textbf{C2 (fidelity).} A fairness edit is worthless if the result no longer
reads as its original driver. Each edit is therefore bounded to a small
neighborhood of the recorded trajectory and scored by a frozen driver-identity
discriminator inside the editing objective (\S\ref{sec:phys-validity}).
\textbf{C3 (the wrong target).} Equal service is the wrong target: a business
district should receive more service than a quiet residential block. The
objective must score only the service variation demand leaves unexplained
(\S\ref{sec:objective}).
\textbf{C4 (level up, not down).} Measured fairness can improve while the
under-served gain nothing, and a demand-only editor does exactly that on our
data (\S\ref{sec:leveling}). Leveling up needs a second channel that adds real
taxi presence where it is missing.
\textbf{C5 (surviving training).} Sparse edits are averaged away by training
that weights every demonstration equally. Their influence must be preserved
during training, and the recovered fairness traced to the edits rather than to
the reweighting itself (\S\ref{sec:downstream}).
